\section{A finite pole lattice and pole-free Lerch closure}
\label{clc:sec:lattice}

Assume now that the positive scales are commensurate: there exists $Q>0$
with $d_j=Q/q_j\in\N$. For rational scales one may choose
\begin{equation}\label{clc:eq:Q}
 Q=\frac{\operatorname{lcm}(\operatorname{num}q_1,\ldots,
                              \operatorname{num}q_r)}
         {\gcd(\operatorname{den}q_1,\ldots,
                              \operatorname{den}q_r)}.
\end{equation}
This is a common multiple of the scales, not a common denominator.
Put
\begin{equation}\label{clc:eq:period-data}
 \sigma=(-1)^{\sum_jd_j},\qquad
 \PP=\bigcup_{j=1}^r\{q_j,2q_j,\ldots,Q\},\qquad
 d_\rho=\#\{j:\rho/q_j\in\mathbb Z\}.
\end{equation}
Then $G_{\bq}(p+Q)=\sigma G_{\bq}(p)$, and every positive pole is a
translate of a member of $\PP$. Define its finite Laurent data by
\begin{equation}\label{clc:eq:Laurent}
 G_{\bq}(\rho+z)=\sum_{k=1}^{d_\rho}a_{\rho,k}z^{-k}+O(1).
\end{equation}

\begin{lemma}[Elementary commensurate kernel]\label{clc:lem:kernel}
For $\mu<0$,
\begin{equation}\label{clc:eq:kernel}
 K_{\bq}(\mu)=
 -\frac{\displaystyle\sum_{\rho\in\PP}\sum_{k=1}^{d_\rho}
 a_{\rho,k}e^{\rho\mu}\frac{\mu^{k-1}}{(k-1)!}}
 {1-\sigma e^{Q\mu}}.
\end{equation}
If $\sigma=1$, the simple Laurent coefficients satisfy
\begin{equation}\label{clc:eq:residue-sum}
 \sum_{\rho\in\PP}a_{\rho,1}=0.
\end{equation}
Formula \eqref{clc:eq:kernel} continues to all real $\mu$, with its
removable value at zero understood when $\sigma=1$.
\end{lemma}
\begin{proof}
Take the inverse transform of $G_{\bq}$ on $0<c<q_*$. For $\mu<0$ close
the contour to the right. Choose the closing vertical lines at a fixed
nonpole residue class modulo $Q$. The trigonometric denominators are
then uniformly separated from zero near the real axis, and give
exponential decay at large imaginary part. The closing vertical
integral tends to zero by $e^{\mu\Rea p}$; the horizontal integrals
vanish as their heights tend to infinity. The closure is clockwise,
so the sign of the residue sum is negative.

At $p=\rho+nQ$, the coefficient of $(p-\rho-nQ)^{-k}$ is
$\sigma^n a_{\rho,k}$. The residue of its product with $e^{p\mu}$ is
$\sigma^na_{\rho,k}e^{(\rho+nQ)\mu}\mu^{k-1}/(k-1)!$.
The convergent geometric sum over $n\ge0$ proves
\eqref{clc:eq:kernel}.

When $\sigma=1$, integrate $G_{\bq}$ around a rectangle whose two
vertical edges differ by $Q$, enclosing one positive period of poles.
The vertical integrals cancel by periodicity, and the horizontal ones
vanish by exponential decay. The residue theorem gives
\eqref{clc:eq:residue-sum}. In \eqref{clc:eq:kernel} the numerator
therefore vanishes at zero, so its possible real singularity is
removable. The original inverse transform is holomorphic in a strip
around the real axis. Analytic continuation from $\mu<0$ along the
real axis proves the last assertion. Possible more distant complex
zeros of the denominator are not being declared removable here.
\end{proof}

\begin{theorem}[Finite Lerch closure at every complex order]
\label{clc:thm:finite}
Let $\Rea A>-q_*$ and $b_\rho=(A+\rho)/Q$. For $\mu<0$ and $W\in\C$,
\begin{align}
 \CC_{\bq,A}(W;\mu)
 =-\sum_{\rho\in\PP}\sum_{k=1}^{d_\rho}a_{\rho,k}e^{\rho\mu}
 \sum_{j=0}^{k-1}
 &\frac{(-1)^j\mu^{k-1-j}(W)_j}
 {j!(k-1-j)!Q^{W+j}}\notag\\[-2pt]
 &\hspace{-24mm}\cdot\Phi(\sigma e^{Q\mu},W+j,b_\rho).
 \label{clc:eq:finite}
\end{align}
The whole expression continues to the real axis. If $\sigma=1$ and
$\mu>0$, its two Lerch boundary values agree. At $\mu=0$ the following
central expression is interpreted in its pole-free form below:
\begin{equation}\label{clc:eq:central-raw}
 \CC_{\bq,A}(W;0)=
 \sum_{\rho\in\PP}\sum_{k=1}^{d_\rho}
 \frac{(-1)^ka_{\rho,k}}{(k-1)!Q^{W+k-1}}
 (W)_{k-1}\Phi(\sigma,W+k-1,b_\rho).
\end{equation}
\end{theorem}
\begin{proof}
For $\Rea W>0$, insert \eqref{clc:eq:kernel} into
\eqref{clc:eq:base-fractional}. Expand
$(\mu-u)^{k-1}/(k-1)!$ in powers of $u$, and use
\[
 \frac1{1-\sigma e^{Q(\mu-u)}}
 =\sum_{n=0}^\infty\sigma^ne^{nQ\mu-nQu}.
\]
The absolute convergence is immediate for $\mu<0$ from
$\Rea(A+\rho)>0$. Each Laplace integral is
$\Gamma(W+j)/(A+\rho+nQ)^{W+j}$.
Dividing by $\Gamma(W)$ produces $(W)_j$, proving
\eqref{clc:eq:finite} in the initial half-plane. Since the Lerch argument
has modulus less than one, the right side is entire in $W$. Theorem
\ref{clc:thm:addition} proves equality everywhere.

Initially with $\Rea W$ sufficiently large, taking $\mu\to0^-$ is
justified in the summed integrals and leaves only $j=k-1$. This gives
\eqref{clc:eq:central-raw}. Its extension to all orders follows from
the pole-free identities in Corollary~\ref{clc:cor:polefree}.

For the real continuation when $\sigma=1$, continue the equality from
negative $\mu$ through the upper and lower parts of the holomorphic
strip. The Lerch jump at
$z=e^{Q\mu}>1$ is, up to the common boundary-orientation sign,
\[
 \frac{2\pi i}{\Gamma(W+j)}
 e^{-b_\rho Q\mu}(Q\mu)^{W+j-1}.
\]
For $\Rea(W+j)>0$, this jump follows by indenting the Lerch Euler
integral at its simple pole $u=Q\mu$; its residue is precisely the
displayed coefficient. The general order follows by continuation.
After multiplication by the factors in \eqref{clc:eq:finite}, the
$\rho$ and $j$ dependence in the exponential and in $Q$ cancels.
For fixed $k\ge2$, the remaining coefficient is proportional to
\[
 \sum_{j=0}^{k-1}\frac{(-1)^j}{j!(k-1-j)!}=0.
\]
For $k=1$, the total remaining coefficient is proportional to
$\sum_\rho a_{\rho,1}=0$. This initially proves zero jump away from
the Gamma exceptional points and then everywhere by order continuation.
Thus both boundary values coincide with the analytic ordinary
convolution. For $\sigma=-1$ there is no positive-real Lerch cut to cross.
\end{proof}

\begin{corollary}[Entire central normalization]\label{clc:cor:polefree}
If $\sigma=-1$, every summand of \eqref{clc:eq:central-raw} is entire,
with $\Phi(-1,v,b)=\eta(v,b)$. If $\sigma=1$, choose any fixed reference
$b_*$ with positive real part. Replace the combined $k=1$ terms by
\begin{equation}\label{clc:eq:simple-polefree}
 -Q^{-W}\sum_{\rho\in\PP}a_{\rho,1}
 \bigl[\zeta(W,b_\rho)-\zeta(W,b_*)\bigr].
\end{equation}
For every $k\ge2$ make the replacement
\begin{equation}\label{clc:eq:high-polefree}
 (W)_{k-1}\zeta(W+k-1,b_\rho)
 =(W)_{k-2}\EE_{b_\rho}(W+k-2).
\end{equation}
All these new summands are entire in $W$ and their sum is the desired
central convolution.
\end{corollary}
\begin{proof}
The residues of the two Hurwitz functions in a difference cancel. At
$W=1$ its value is $\psi(b_*)-\psi(b_\rho)$ by
\eqref{clc:eq:stieltjes-convention}. Equation
\eqref{clc:eq:residue-sum} makes subtraction of the reference function
legitimate. For $k\ge2$, the last factor of $(W)_{k-1}$ is
$W+k-2$, giving \eqref{clc:eq:high-polefree}. The function $\EE$ is
entire by its removable normalization. These observations prove the
claim, including every nonpositive-integer resonance.
\end{proof}

The number of terms is finite independently of the orders. It depends
only on the scales and pole multiplicities. No search for an integer
relation among numerical special values enters the proof.
